Orbiter can plot functions using a built-in function tracker.
The functions must be continuous apart from a finite number of poles.
The function can have multiple components, each described using an expression. 
Each expression is specified in Reverse Polish Notation (RPN).
Consider an example. 
A Lissajous curve is defined using coordinate functions of the form 
$$
x = r \sin\big(at + c\big), \quad
y = r \sin\big(bt\big), \quad a,b,c,r \in \bbR.
$$
The terms 
$$
r \sin\big(at + c\big), \quad r \sin\big(bt\big)
$$
are the expressions of the two coordinate functions.
RPN means that the operator is listed after the operands. 
A stack data structure is used to hold temporary values. 
Operators are pushed to the top of the stack using the push commands.
A binary operator pops the two elements from the stack, performs the operation, 
and pushes the resulting value back onto the stack. 
For a unary operator, only one element is popped and replaced by the result.
Here are some examples of expressions rewritten in RPN:
\begin{align*}
\sin(x) &\mapsto \; \mbox{push} \; x \; \sin,\\
a+b & \mapsto \; \mbox{push} \; a \; \mbox{push} \; b \; \mbox{add},\\
a\cdot b & \mapsto \; \mbox{push} \; a \; \mbox{push} \; b \; \mbox{mult}.
\end{align*}
The coordinate functions are enclosed between \verb'-code' and \verb'-code_end' commands. 
Each coordinate function is described in RPN and terminated using a \verb'return' keyword.
By the time the \verb'return' keyword is reached, the RPN expression must have exactly 
one value on the stack which is considered the value of the expression.
Constants are declared between the \verb'-const' and \verb'-const_end' keywords. 
Likewise, variables are declared between the \verb'-var' and \verb'-var_end' keywords.
Picking $a=3$, $b=2,$ $c = \pi/2$ and $r=7$, the function is computed using 

\sourcecode{lissajous}

The sequence
$$
\verb'push t push a mult push c add sin push r mult'
$$
is $r\sin(at+c)$ expressed in RPN.
The constants are defined in the line 
$$
\verb'-const a 3 b 2 c 1.57 r 7 -const_end'
$$
The input variable is defined using the line 
$$
\verb'-var t -var_end'
$$
The sequence 
$$
\verb'-smooth_curve "lissajous" 0.07 2000 15 0 18.85'
$$
defines the name of the output file, the fact that two consecutive 
points are never further than $\epsilon =0.07$ away, 
the fact that points that are 15 or more away from the origin should be ignored, 
and the fact that the variable $t$ loops over the range $[0,18.85]$ 
with a default of $2000$ steps.
The evaluator automatically reduces the step-size 
if consecutive image points are more than $\epsilon$ apart.  
The code to produce the plot is

\sourcecode{lissajous_plot}

The plot is shown below
\orbiteroutput{GRAPHICS/WRAPPERS/lissajous_0_000_draw}



\bigskip


We can turn it into a 3D plot by using the $t$ value for the $z$ coordinate. 
The 
function is computed using the command 

\sourcecode{lissajous_3d}

The code to produce the 3D plot is

\sourcecode{lissajous_3d_plot}

The 3D curve is shown below
\orbiteroutput{GRAPHICS/WRAPPERS/lissajous_3d_draw}

